\section{Technical details for the variational and finite--$N$ arguments}

This appendix collects supplementary arguments used at the delicate points of the radial variational problem and of the finite--$Z$ propagation.  The purpose is to record the functional--analytic justification of the compactness and first--variation steps, the support and uniqueness issues in the global maximization argument, and the explicit estimates entering the numerical coefficients of \Cref{thm:finite-Z}.

\subsection{Admissible-class reduction}

The HPS definition uses radial probability measures belonging to $H^{-1}(\R^3)$ and with finite second moment.  For the lower bound we enlarge the class after radial reduction.  The lower bound is therefore proved on a larger set, while the explicit optimizer \eqref{eq:rho-star} is an $H^{-1}$ probability density with compact support bounded away from the origin.  Hence the lower and upper bounds meet inside the original HPS admissible class.  The radial reduction of HPS identifies the value of the full variational infimum with the radial infimum; the uniqueness result proved here is, however, a uniqueness statement within the radial class.

\subsection{Compactness under the exponential moment constraints}

Weak compactness alone does not preserve the constraints
\[
 A[\sigma]=B[\sigma]=1,
\]
because exponentially small masses can escape to large $|t|$ while carrying a nonvanishing exponential moment.  The proof of \Cref{lem:existence} therefore uses the value of the variational functional in addition to tightness.  A normalized two--point measure gives an explicit competitor with $J>0$, so the supremum $M$ is strictly positive.  Tightness and bounded total mass yield narrow convergence along a maximizing subsequence, and the corresponding product measures converge narrowly as well.  Since $k(t-u)$ is bounded and continuous, this gives
\[
 J[\sigma_n]\longrightarrow J[\sigma]=M.
\]
If the weak limit satisfied $A[\sigma]B[\sigma]<1$, amplitude and translation normalization would multiply $J$ by $(A[\sigma]B[\sigma])^{-1}>1$, producing a normalized competitor with value larger than $M$.  Consequently the exponential moments are preserved at a maximizer.

\subsection{First variation and the obstacle relation}

The first variation is taken directly in the scale-- and translation--invariant quotient
\[
 \frac{J[\sigma]}{A[\sigma]B[\sigma]}.
\]
Adding an atom $\varepsilon\delta_t$ gives the global obstacle inequality.  Equality on the support follows from
\[
 \int\left[M(\ee^t+\ee^{-t})-2(k*\sigma)(t)\right]\dd\sigma(t)=0
\]
together with nonnegativity of the integrand.  This yields the Euler--Lagrange relation used in the main text without requiring a separate infinite-dimensional constraint-qualification argument.

\subsection{Absence of atoms and endpoint regularity}

The cusp of $k$ at the origin is essential.  An atom of mass $m$ in $\sigma$ produces an upward derivative jump $2m$ in $k*\sigma$ and hence a downward jump $-4m$ in the derivative of the nonnegative obstacle.  Such a jump is incompatible with contact at a local minimum, and therefore a maximizing measure is atomless.

For a compactly supported atomless finite measure, convolution with the Lipschitz kernel $k$ is $C^1$.  The only discontinuity of the one-sided derivatives of $k$ occurs at $u=t$, which carries zero mass, and continuity follows by dominated convergence after separating a small neighborhood of $t$.  Hence the obstacle is $C^1$.  At each extremal point of the support both its value and its first derivative vanish, which justifies the smooth-pasting identities used in \Cref{lem:endpoints}.

\subsection{Conditional negativity for signed measures}

Let $\eta$ be a finite signed measure satisfying the two exponential moment constraints.  The function
\[
 p(t)=\int\ee^{|t-u|}\dd\eta(u)
\]
vanishes outside the candidate interval and belongs to $H_0^1$.  Distributionally,
\[
 (\partial_t^2-1)p=2\eta,
\]
and hence $\eta\in H^{-1}(\R)$.  Fourier transformation gives
\[
 \widehat\eta(\xi)=-\frac{1+\xi^2}{2}\widehat p(\xi).
\]
Together with
\[
 \widehat{k}(\xi)=\frac{2(2-\xi^2)}{(1+\xi^2)(4+\xi^2)},
\]
this yields
\[
 \widehat{k}(\xi)|\widehat\eta(\xi)|^2
 =\frac12\left(-\xi^2+5-\frac{18}{\xi^2+4}\right)|\widehat p(\xi)|^2,
\]
and therefore the identity \eqref{eq:Jeta-identity}.  The finite-measure pairing follows by standard approximation in $H^{-1}$.  The strict sign needed in the proof is then a consequence of
\[
 2a_*<\frac{\pi}{\sqrt5}
\]
and the Dirichlet Poincar\'e inequality on $[-a_*,a_*]$.

\subsection{Support geometry in the global argument}

No connectedness assumption on $\supp\sigma$ is required.  Only its extremal points
\[
 a=\inf\supp\sigma,
 \qquad
 b=\sup\supp\sigma
\]
enter the argument.  Compactness and smooth pasting at these points give the endpoint identities and hence $a+b=0$.  If a maximizer has value at least $M_*$, its extremal support points lie inside $[-a_*,a_*]$, irrespective of possible gaps in the support.  Thus the difference measure $\eta=\sigma-\sigma_*$ is supported in the full candidate interval, exactly as required by \Cref{lem:conditional-negativity}.

\subsection{Radial uniqueness and its scope}

The proof of \Cref{thm:main} establishes uniqueness of the normalized maximizer of the logarithmic radial problem and therefore uniqueness, up to dilation, of the radial minimizer.  HPS radialization at $s=3$ proves equality between the full and radial infima, but the radialization inequality is not used here in a strict form.  The result therefore does not imply uniqueness of minimizers in the full nonradial HPS admissible class.

\subsection{Propagation of the exact constant through the HPS finite--$N$ estimates}

The finite--$Z$ theorem uses the symbolic cubic-weight estimates in \cite[Sec.~7.3]{HPS2025}.  In particular, the finite--configuration estimate has the form
\[
 (1-r^2)N\beta_3-\frac1r
 \le
 \left(1+\frac{r^2}{5}+\frac{r^3}{35}\right)\alpha_{N,3}(N-1),
\]
while the weighted kinetic estimate bounds the right--hand side by the same factor times
\[
 Z+c_0ZN^{-2/3}.
\]
Thus the mean--field constant $\beta_3$ remains symbolic throughout the finite--$N$ chain.  The exact value obtained in \Cref{thm:main} can therefore be inserted before the final numerical optimization, with no modification of the underlying many--body inequalities.

\subsection{Endpoint optimization for the $Z^{1/3}$ coefficient}

After isolating the lower--order coefficients, the contradiction argument reduces to bounding
\[
 F_{\beta_3}(x)
 =
 3\left(\frac3{10}\right)^{1/3}\beta_3^{-2/3}x^{1/3}
 +c_0\beta_3^{-1}x^{-2/3},
 \qquad x=\frac NZ.
\]
Its derivative has the form
\[
 F'_{\beta_3}(x)
 =\frac{x^{-5/3}}{3}(Ax-2B),
 \qquad A,B>0,
\]
so $F_{\beta_3}$ has a unique interior minimum and no interior maximum.  It is therefore sufficient to compare the endpoints of the admissible interval.

For $Z\ge4$, Lieb's bound gives
\[
 x<2+\frac1Z\le\frac94.
\]
Hence
\[
 \beta_3^{-1}\le x\le\frac94.
\]
Evaluation of the two endpoints gives
\[
 F_{\beta_3}(\beta_3^{-1})
 =3.81147042506\ldots,
 \qquad
 F_{\beta_3}(9/4)
 =3.78401155302\ldots.
\]
Thus the left endpoint is the maximum on the interval relevant to \Cref{thm:finite-Z}, and the rounded coefficient $3.812$ is valid.  Retaining the sharper upper endpoint $9/4$ avoids the unnecessary loss that would arise from enlarging the interval to $[\beta_3^{-1},5/2]$.

\subsection{Certified evaluation of the explicit constants}

The decimal constants in \Cref{thm:finite-Z} are chosen on the safe side of explicit closed expressions.  Interval evaluation gives
\[
 \beta_3^{-1}=1.08632517518\ldots<1.086326,
\]
\[
 \bar a_2=0.01293244256\ldots<0.01294,\qquad
 \bar a_3=0.18184084198\ldots<0.18185,
\]
and
\[
 \bar a_4=0.01940243839\ldots<0.01941,
 \qquad
 \sup F_{\beta_3}=3.81147042506\ldots<3.812.
\]
All quantities are obtained from elementary closed formulas involving $\exp$, $\arctan$, radicals and $\pi$; no numerical optimization is used in the theorem.

\subsection{Limitations of the present argument}

The result sharpens the finite--$Z$ estimate within the cubic HPS mean--field framework but does not prove the ionization conjecture
\[
 N_c(Z)\le Z+C.
\]
Nor does it establish that $\beta_3^{-1}$ is the optimal coefficient among all weighted multiplier approaches.  Further improvement must exploit information absent from the scalar radial mean--field problem, such as additional positive terms in the multiplier identity or genuinely fermionic occupation constraints.
